\documentclass{article}

\title{Ch. 24 HW \#1 Solutions}
\author{Principles of Macroeconomics}
\date{\today}

\begin{document}

\maketitle

\section*{Questions for Review}

\begin{enumerate}

\item Individuals generate income by selling goods and services. Every time someone sells a good or service, there's a buyer on the other side of that transaction. That buyer is spending the same amount of money that the seller is making. That means every dollar of income earned is also a dollar of expenditure. 

\item The production of a luxury car car contributes more to GDP because it has a greater market value.

\item The total contribution to GDP is \$3 because the bread is the only final good bought or sold.

\item It does not because the goods weren't produced in the current period.

\end{enumerate}

\section*{Problems and Applications}

\begin{enumerate}

\setcounter{enumi}{8}

\item

	\begin{enumerate}
	
	\item GDP is \$180 because the only final good produced in the economy is the bread.
	
	\item Farmer: $100 - 0 = 100$ \\
		Miller: $150 - 100 = 50$ \\
		Baker: $180 - 150 = 30$
		
	\item The total value added is $100 + 50 + 30 = 180$ which is the same as GDP. This suggests a third way of calculating GDP: summing every producer's value added. 
	
	\end{enumerate}

\end{enumerate}

\end{document}